% Generated draft: requires mathematical and source review.
\documentclass[11pt]{article}
\usepackage[margin=25mm]{geometry}
\usepackage{fontspec}
\setmainfont{Latin Modern Roman}
\usepackage{amsmath,amssymb,mathtools,unicode-math}
\setmathfont{Latin Modern Math}
\usepackage{enumitem,needspace,xurl,hyperref}
\setlength{\parindent}{0pt}
\setlength{\parskip}{5pt}
\setlength{\emergencystretch}{3em}
\newcommand{\N}{\mathbb{N}}
\newcommand{\Z}{\mathbb{Z}}
\newcommand{\Q}{\mathbb{Q}}
\newcommand{\R}{\mathbb{R}}
\newcommand{\C}{\mathbb{C}}
\newcommand{\F}{\mathbb{F}}
\newcommand{\A}{\mathbb{A}}
\newcommand{\PP}{\mathbb{P}}
\newcommand{\E}{\mathbb{E}}
\newcommand{\eps}{\varepsilon}
\newcommand{\dd}{\,\mathrm{d}}
\begin{document}
\section*{UnsolvedMath ID: 30002694}
\textbf{Source problem code:} catalogue entry
\subsection{Definitions and mathematical statement}
Let \((M,g)\) be a smooth Riemannian manifold with Levi--Civita connection \(\nabla\) and curvature tensor \(R\). The metric is homogeneous if its isometry group acts transitively on \(M\), and Einstein if \(\operatorname{Ric}_g=\lambda g\) for a constant \(\lambda\). A geodesic-orbit metric is one for which every geodesic is an orbit of a one-parameter group of isometries. A homogeneous presentation \(M=G/H\) is naturally reductive if it admits an \(\operatorname{Ad}(H)\)-invariant decomposition \(\mathfrak g=\mathfrak h\oplus\mathfrak m\) such that the metric inner product on \(\mathfrak m\) satisfies
\[
\langle [X,Y]_{\mathfrak m},Z\rangle+\langle Y,[X,Z]_{\mathfrak m}\rangle=0
\qquad(X,Y,Z\in\mathfrak m).
\]
Here natural reductivity and the geodesic-orbit property are understood to hold for a suitable transitive isometry group; the supplied record does not specify a particular group for the general questions.

The source denotes by \(k_o\) the constant Jacobi osculating rank and by \(k_S\) the Singer invariant. It defines
\[
 k_t(M,g)=\min\bigl\{k\in\Z_{\geq0}:Y_0^{(k+3,2)}R\equiv0\bigr\},
\]
where \(Y_0^{(k+3,2)}R\) is called the completely traceless Cartan component. For expressing the finiteness questions, set \(k_t=+\infty\) when this set is empty. The supplied record does not give the construction of this component, the precise definition and constancy requirement for \(k_o\), or the indexing convention for \(k_S\). These source-dependent definitions must be recovered before the numerical questions and inequalities constitute a fully specified mathematical problem; no convention for them is silently imposed here.

Subject to those unresolved definitions, the supplied problem asks the following.
\begin{enumerate}
\item Determine \(k_t\) for the intended homogeneous metrics on
\[
 \operatorname{Sp}(2)/\operatorname{SU}(2),\qquad
 \bigl(\operatorname{SO}(3)\times\operatorname{SU}(3)\bigr)/U^\bullet(2),\qquad
 \operatorname{SU}(3)/T^2,
\]
and on Kaplan's six-dimensional geodesic-orbit nilmanifold. The subgroup embeddings, the meaning of \(U^\bullet(2)\), and the intended metrics are not specified in the supplied record. The nilmanifold is named but its Lie algebra and metric are not supplied. Thus these names alone do not specify the requested numerical computations.
\item Compute \(k_o\), \(k_t\), and \(k_S\) for further homogeneous Riemannian spaces. This is an open-ended request; no further list is supplied.
\item For each of the three classes consisting of naturally reductive, geodesic-orbit, and arbitrary homogeneous Riemannian manifolds, determine whether
\[
 \text{for every }(M,g)\text{ in the class, there exists }k\in\Z_{\geq0}
 \text{ such that }Y_0^{(k+3,2)}R\equiv0.
\]
\item Determine the intended converse relating finite \(k_t\) to homogeneity. The clean statement asks whether finite \(k_t\) characterizes complete homogeneous spaces, whereas the original statement merely asks whether the converse holds. In particular, it is unresolved whether the intended assertion is
\[
 \text{for every complete Riemannian }(M,g),\qquad
 k_t(M,g)<\infty\ \Longleftrightarrow\ (M,g)\text{ is homogeneous},
\]
or whether finiteness itself is intended to imply both completeness and homogeneity. These alternatives are recorded without selecting one. The use of \(k_t\) on potentially nonhomogeneous manifolds also requires the missing tensor definition.
\item Determine whether every homogeneous Einstein Riemannian manifold has \(k_o<\infty\), with \(k_o\) understood in the source's intended sense.
\item Characterize the equality cases in the source-stated inequalities
\[
 k_t\leq k_o,\qquad k_S\leq k_o,
\]
and, in the Einstein case,
\[
 k_S\leq k_t.
\]
The record presents these as inequalities whose equality cases are sought. It does not supply their proofs or full applicability conditions, including any finiteness or constancy hypotheses; those conditions require review.
\end{enumerate}
This is a statement-only transcription with explicit unresolved qualifications, not a verification of the inequalities, the invariant definitions, or any literature-status claim.
\subsection{Short English statement}
The problem asks for curvature-derivative invariants of specified and further homogeneous Riemannian spaces, universal finiteness results, a converse relating finite \(k_t\) to complete homogeneity, and equality cases in several stated inequalities. Essential invariant conventions, example metrics and subgroup embeddings, and the precise converse remain unspecified. The supplied registry labels the problem partially solved; this draft does not independently verify that assessment.
\subsection{Sources}
\begin{itemize}
\item Local registry: \texttt{data/unsolvedmath/problems.json}, with \texttt{catalog/problems.json} as the supplementary registry; supplied complete record for permanent UnsolvedMath ID \(30002694\), titled ``Curvature-Derivative Invariants of Homogeneous Spaces,'' exact source problem code \texttt{catalogue entry}. Registry collection: Oberwolfach Reports Open Problems. No individual proposer is supplied; the proposed year is recorded as 2014.
\item Supplied report citation: \emph{Mini-Workshop: Einstein Metrics, Ricci Solitons and Ricci Flow under Symmetry Assumptions} (2015), Oberwolfach Reports, pp.\ 2529--2568, DOI \texttt{10.4171/owr/2014/45}. \url{https://doi.org/10.4171/owr/2014/45}. The literature-source label is ``Oberwolfach Report 45/2014.'' The supplied citation's year 2015 and the report identifier/year 2014 are retained without resolving their discrepancy.
\item Supplied literature source: \emph{Flag Manifolds with Five Isotropy Summands}. \url{https://doi.org/10.48550/arXiv.1207.2897}. The record gives no authors, theorem numbers, or explicit correspondence between this reference and particular invariant computations.
\item Supplied dated literature curation, checked at \texttt{2026-08-21}: status ``partially solved.'' Its assessment says that several listed invariants have been computed for important homogeneous examples and special nearly K\"ahler spaces, while the universal finiteness, converse, and equality questions remain open. This is supplied dated curation, not an independently verified literature assessment. The background expressly requests independent verification of time-sensitive status and cited claims.
\item The registry marks the statement \texttt{corrected\_verified} and reports that it was checked against the official Oberwolfach report. This is a supplied provenance claim; no source visit or independent mathematical verification is asserted here.
\end{itemize}
\end{document}
